% ============================================================
%  Chapter 16
% ============================================================

\chapter{Projection Engines}

\chapterepigraph{%
  The same latent structure\\
  can be heard as music,\\
  seen as image,\\
  read as text.\\
  The engine is the same; the output is the choice.
}{Flyxion, \textit{Semantic Infrastructure}}

\sidenote{App.\,B for the\\formal projection\\theory; App.\,N\\for universal\\approximation}

Chapter~15 showed how constraints reduce the search space for inference.
This chapter examines the output side: how shared latent structures project
into text, music, images, and behavior — and why the same underlying
constraint geometry can appear in radically different surface forms.

\section{Outputs as Projections}

In the RSVP framework, any observable output is a projection of the
underlying field configuration $(\Phi, \mathbf{v}, S)$ onto an
observable space. Different output modalities correspond to different
projections.

\begin{definition}{Output Projection}{output-proj}
  An \emph{output projection} is a map $\Pi_\text{out} : \mathcal{F} \to Y$
  from the field configuration space $\mathcal{F}$ to an observable
  output space $Y$. Different projection maps $\Pi_\text{out}$ define
  different output modalities.
\end{definition}

The key insight is that $\mathcal{F}$ — the space of field configurations —
is shared across all modalities. The differences between text, music,
and images are differences in the projection, not differences in the
underlying structure.

\section{Text as Projection}

Language is perhaps the most compressed of all projections. A sentence
of $n$ words drawn from a vocabulary $V$ occupies a point in the discrete
space $V^n$. The sentence's meaning, however, lives in the continuous
semantic manifold $M_s$ of Chapter~23 (Appendix~L).

The text projection $\Pi_\text{text} : M_s \to V^n$ maps semantic
configurations to word sequences. This projection is enormously many-to-one:
the fiber $\Pi_\text{text}^{-1}(w)$ over a sentence $w$ contains the
entire equivalence class of semantic configurations that would be expressed
by $w$. Synonymy, paraphrase, and translation are all manifestations of
non-trivial fibers in the text projection.

\section{Music as Projection}

Music projects the same latent structure into a different observable space:
time-frequency energy distributions, spectral sequences, rhythmic patterns.

\begin{definition}{Musical Projection}{music-proj}
  The \emph{musical projection} is $\Pi_\text{music} : M_s \to \mathcal{M}_\text{mus}$
  where $\mathcal{M}_\text{mus}$ is the space of musical performances
  (continuous waveforms, or symbolic scores, depending on level of analysis).
\end{definition}

The shared latent structure between text and music is most visible in
their common organizational principles: narrative arc, tension and release,
repetition and variation, hierarchical phrase structure. These are not
metaphors — they are manifestations of the same accessibility geometry
projected through different modalities.

\section{Images as Projection}

Visual images project the latent structure into two-dimensional spatial
patterns. The projection $\Pi_\text{image} : \mathcal{F} \to \mathcal{I}$
(where $\mathcal{I}$ is the space of images) has an interesting property:
it tends to be more faithful than text or music projections, in the sense
that natural images often have smaller fibers — there are fewer semantic
configurations consistent with a given image than there are configurations
consistent with a given sentence.

This is why ``a picture is worth a thousand words'': the image projection
retains more information than the text projection.

\section{Shared Latent Structures}

The deepest claim of the projection engine framework is that text,
music, images, and behavior share a common latent space.

\begin{namedthm}{The Projection Engine Principle}
  Multiple observable modalities are different projections of the same
  underlying accessibility geometry. Translation between modalities is
  not conversion between different objects — it is comparison of different
  projections of the same object. The universal structure is not the
  text, not the music, not the image, but the field configuration
  $(\Phi, \mathbf{v}, S)$ from which all are projected.
\end{namedthm}

\begin{exercise}
  Identify three pairs of text and image that are ``projections of the
  same thing'' — where the text and image share the same latent structure.
  For each pair, describe what the shared latent structure is.
  What information does each modality contain that the other lacks?
\end{exercise}

\begin{exercise}
  Music is sometimes described as ``temporal'' and images as ``spatial.''
  In the projection engine framework, is this distinction fundamental or
  an artifact of the projection? Construct an argument that both modalities
  project structures that are neither purely temporal nor purely spatial.
\end{exercise}

\begin{exercise}[title={Technical}]
  Universal approximation theorems (Appendix~N) say that neural networks
  can approximate any continuous function. Interpret this in terms of
  projection engines: what is being approximated, and what is the
  projection? What does ``sufficient depth'' mean in terms of the
  accessibility landscape?
\end{exercise}

